% --- Archon physics formalization source begin ---
% archon:physics
% archon:covers PhyXMiniProblems/problem_phyx_mini_0139.lean
% archon:source-report reports/phyx_mini/problem_phyx_mini_0139.source.json

\chapter{Physics problem phyx\_mini\_0139}
\label{ch:PhyXMiniProblems_problem_phyx_mini_0139}

\paragraph{Problem source.}
\#\# Physical scenario

An external rearview car mirror is convex with a radius of curvature of \textbackslash{}(16.0\textbackslash{},\textbackslash{}mathrm\{m\}\textbackslash{}) (figure). An object is \textbackslash{}(10.0\textbackslash{},\textbackslash{}mathrm\{m\}\textbackslash{}) from the mirror.

\#\# Question

Determine the image's magnification for the object.

\#\# Answer choices

- A: A: +0.24

- B: B: +0.34

- C: C: +0.44

- D: D: +0.54

\#\# Figure caption (auxiliary; use the image as primary evidence)

The image shows the reflective view in a car's side mirror. In the mirror, there is a highway scene where a large red truck is approaching from behind. The landscape around the road features tall green trees and a clear blue sky with white clouds. The edge of the car mirror and part of the car door are visible, indicating that the image is taken from inside the vehicle. There is no text present in the image.

\#\# Dataset metadata

Geometrical Optics; Physical Model Grounding Reasoning, Multi-Formula Reasoning

\paragraph{Recorded answer/context.}
C: C: +0.44

\paragraph{Figure/image path.}
/root/proposal\_for\_physic/science-mango/phyx\_mini\_run/phyx\_data/test\_image/139.png

\paragraph{Formalization target.}
create a compiling Lean file with sorry bodies at `PhyXMiniProblems/problem\_phyx\_mini\_0139.lean`. The Lean declarations must preserve the physical quantities, dimensions or dimensional roles, figure labels, governing-law hypotheses, and final relation expressed by this problem.
Use Mathlib/Physlib names found through LeanExplore where available. If a physics API is missing, introduce faithful local abstractions rather than scalar placeholder aliases.

\begin{theorem}[Physics formalization target]
\label{thm:physics:phyx_mini_0139:target}
\lean{PhyXMiniProblems.ProblemPhyXMini0139.problem_phyx_mini_0139}
\uses{def:physics:phyx-mini-0139:phyxminiproblems-problemphyxmini0139-convexrearviewmirrorsetup, def:physics:phyx-mini-0139:phyxminiproblems-problemphyxmini0139-matchesproblemstatement, def:physics:phyx-mini-0139:phyxminiproblems-problemphyxmini0139-matchesprimaryfigure, def:physics:phyx-mini-0139:phyxminiproblems-problemphyxmini0139-usescartesianconvexmirrorconvention, def:physics:phyx-mini-0139:phyxminiproblems-problemphyxmini0139-satisfiesconvexmirrorfocallaw, def:physics:phyx-mini-0139:phyxminiproblems-problemphyxmini0139-satisfiesgaussianmirrorequation, def:physics:phyx-mini-0139:phyxminiproblems-problemphyxmini0139-satisfiestransversemagnificationlaw, def:physics:phyx-mini-0139:phyxminiproblems-problemphyxmini0139-answerchoice, def:physics:phyx-mini-0139:phyxminiproblems-problemphyxmini0139-isnearesthundredthreadout}
The assigned autoformalize agent should translate this physics problem into Lean declarations in the covered file, with theorem and lemma proof bodies written as `by sorry`.
\end{theorem}
\begin{proof}
This is an autoformalization task, not a proof task. Produce statements that can later be proved by the physics prover without weakening the physical model.
\end{proof}

% --- Archon declaration topology begin ---
\section{Declaration topology}
\begin{definition}[Optical Length]
  \label{def:physics:phyx-mini-0139:phyxminiproblems-problemphyxmini0139-opticallength}
  \lean{PhyXMiniProblems.ProblemPhyXMini0139.OpticalLength}
  A signed optical length whose scalar readout varies coherently with units
\end{definition}
\begin{proof}
  This is the declared model definition of Optical Length.
\end{proof}
\begin{definition}[length In Meters]
  \label{def:physics:phyx-mini-0139:phyxminiproblems-problemphyxmini0139-lengthinmeters}
  \lean{PhyXMiniProblems.ProblemPhyXMini0139.lengthInMeters}
  \uses{def:physics:phyx-mini-0139:phyxminiproblems-problemphyxmini0139-opticallength}
  The signed scalar readout of an optical length in SI metres
\end{definition}
\begin{proof}
  This is the declared model definition of length In Meters.
\end{proof}
\begin{definition}[Spherical Mirror Kind]
  \label{def:physics:phyx-mini-0139:phyxminiproblems-problemphyxmini0139-sphericalmirrorkind}
  \lean{PhyXMiniProblems.ProblemPhyXMini0139.SphericalMirrorKind}
  The two spherical-mirror kinds, viewed from the incident-light side
\end{definition}
\begin{proof}
  This is the declared model definition of Spherical Mirror Kind.
\end{proof}
\begin{definition}[Mirror Role]
  \label{def:physics:phyx-mini-0139:phyxminiproblems-problemphyxmini0139-mirrorrole}
  \lean{PhyXMiniProblems.ProblemPhyXMini0139.MirrorRole}
  The use made of the mirror in the physical scenario
\end{definition}
\begin{proof}
  This is the declared model definition of Mirror Role.
\end{proof}
\begin{definition}[Figure Element]
  \label{def:physics:phyx-mini-0139:phyxminiproblems-problemphyxmini0139-figureelement}
  \lean{PhyXMiniProblems.ProblemPhyXMini0139.FigureElement}
  Distinct physical elements recognizable in the supplied photograph
\end{definition}
\begin{proof}
  This is the declared model definition of Figure Element.
\end{proof}
\begin{definition}[Convex Rearview Mirror Setup]
  \label{def:physics:phyx-mini-0139:phyxminiproblems-problemphyxmini0139-convexrearviewmirrorsetup}
  \lean{PhyXMiniProblems.ProblemPhyXMini0139.ConvexRearviewMirrorSetup}
  \uses{def:physics:phyx-mini-0139:phyxminiproblems-problemphyxmini0139-opticallength, def:physics:phyx-mini-0139:phyxminiproblems-problemphyxmini0139-sphericalmirrorkind, def:physics:phyx-mini-0139:phyxminiproblems-problemphyxmini0139-mirrorrole, def:physics:phyx-mini-0139:phyxminiproblems-problemphyxmini0139-figureelement}
  The physical mirror configuration. radiusOfCurvatureMagnitude and objectDistance are positive magnitudes. signedFocalLength and signedImageDistance use the Cartesian convention in which a convex mirror has negative focal length and its virtual image lies behind the mirror at negative image distance. No value is assigned here to the requested transverseMagnification
\end{definition}
\begin{proof}
  This is the declared model definition of Convex Rearview Mirror Setup.
\end{proof}
\begin{definition}[Matches Problem Statement]
  \label{def:physics:phyx-mini-0139:phyxminiproblems-problemphyxmini0139-matchesproblemstatement}
  \lean{PhyXMiniProblems.ProblemPhyXMini0139.MatchesProblemStatement}
  \uses{def:physics:phyx-mini-0139:phyxminiproblems-problemphyxmini0139-lengthinmeters, def:physics:phyx-mini-0139:phyxminiproblems-problemphyxmini0139-convexrearviewmirrorsetup}
  The numerical and categorical data stated in the problem: an external convex rearview mirror of curvature-radius magnitude 16.0 m, imaging the truck at an object distance of 10.0 m. Neither the signed image distance nor the magnification is fixed here
\end{definition}
\begin{proof}
  This is the declared model definition of Matches Problem Statement.
\end{proof}
\begin{definition}[Matches Primary Figure]
  \label{def:physics:phyx-mini-0139:phyxminiproblems-problemphyxmini0139-matchesprimaryfigure}
  \lean{PhyXMiniProblems.ProblemPhyXMini0139.MatchesPrimaryFigure}
  \uses{def:physics:phyx-mini-0139:phyxminiproblems-problemphyxmini0139-convexrearviewmirrorsetup}
  Qualitative readouts from the primary photograph. The truck, highway, trees, and sky occur in the mirror's reflected view; the mirror frame and host car body are directly visible. The image contains no text or metric annotations
\end{definition}
\begin{proof}
  This is the declared model definition of Matches Primary Figure.
\end{proof}
\begin{definition}[Uses Cartesian Convex Mirror Convention]
  \label{def:physics:phyx-mini-0139:phyxminiproblems-problemphyxmini0139-usescartesianconvexmirrorconvention}
  \lean{PhyXMiniProblems.ProblemPhyXMini0139.UsesCartesianConvexMirrorConvention}
  \uses{def:physics:phyx-mini-0139:phyxminiproblems-problemphyxmini0139-lengthinmeters, def:physics:phyx-mini-0139:phyxminiproblems-problemphyxmini0139-convexrearviewmirrorsetup}
  Physical signs for the chosen Cartesian mirror convention. These are qualitative branch conditions, not numerical values of the requested magnification or of the derived image distance
\end{definition}
\begin{proof}
  This is the declared model definition of Uses Cartesian Convex Mirror Convention.
\end{proof}
\begin{definition}[Satisfies Convex Mirror Focal Law]
  \label{def:physics:phyx-mini-0139:phyxminiproblems-problemphyxmini0139-satisfiesconvexmirrorfocallaw}
  \lean{PhyXMiniProblems.ProblemPhyXMini0139.SatisfiesConvexMirrorFocalLaw}
  \uses{def:physics:phyx-mini-0139:phyxminiproblems-problemphyxmini0139-convexrearviewmirrorsetup}
  For a paraxial convex spherical mirror, the signed focal length is minus half the positive radius-of-curvature magnitude. Requiring the equality in every unit system makes it independent of the metre readout used for the data
\end{definition}
\begin{proof}
  This is the declared model definition of Satisfies Convex Mirror Focal Law.
\end{proof}
\begin{definition}[Satisfies Gaussian Mirror Equation]
  \label{def:physics:phyx-mini-0139:phyxminiproblems-problemphyxmini0139-satisfiesgaussianmirrorequation}
  \lean{PhyXMiniProblems.ProblemPhyXMini0139.SatisfiesGaussianMirrorEquation}
  \uses{def:physics:phyx-mini-0139:phyxminiproblems-problemphyxmini0139-convexrearviewmirrorsetup}
  The signed Gaussian spherical-mirror equation 1/f = 1/dₒ + 1/dᵢ, written in the division-free, dimensionally homogeneous form f * (dₒ + dᵢ) = dₒ * dᵢ in every unit system
\end{definition}
\begin{proof}
  This is the declared model definition of Satisfies Gaussian Mirror Equation.
\end{proof}
\begin{definition}[Satisfies Transverse Magnification Law]
  \label{def:physics:phyx-mini-0139:phyxminiproblems-problemphyxmini0139-satisfiestransversemagnificationlaw}
  \lean{PhyXMiniProblems.ProblemPhyXMini0139.SatisfiesTransverseMagnificationLaw}
  \uses{def:physics:phyx-mini-0139:phyxminiproblems-problemphyxmini0139-convexrearviewmirrorsetup}
  The signed transverse-magnification law m = -dᵢ/dₒ, written without division. The two distances are dimensionful while m is dimensionless
\end{definition}
\begin{proof}
  This is the declared model definition of Satisfies Transverse Magnification Law.
\end{proof}
\begin{definition}[Answer Choice]
  \label{def:physics:phyx-mini-0139:phyxminiproblems-problemphyxmini0139-answerchoice}
  \lean{PhyXMiniProblems.ProblemPhyXMini0139.AnswerChoice}
  Labels of the four multiple-choice answers in the source problem
\end{definition}
\begin{proof}
  This is the declared model definition of Answer Choice.
\end{proof}
\begin{definition}[magnification Readout]
  \label{def:physics:phyx-mini-0139:phyxminiproblems-problemphyxmini0139-answerchoice-magnificationreadout}
  \lean{PhyXMiniProblems.ProblemPhyXMini0139.AnswerChoice.magnificationReadout}
  \uses{def:physics:phyx-mini-0139:phyxminiproblems-problemphyxmini0139-answerchoice}
  The dimensionless decimal magnification printed beside each choice
\end{definition}
\begin{proof}
  This is the declared model definition of magnification Readout.
\end{proof}
\begin{definition}[Is Nearest Hundredth Readout]
  \label{def:physics:phyx-mini-0139:phyxminiproblems-problemphyxmini0139-isnearesthundredthreadout}
  \lean{PhyXMiniProblems.ProblemPhyXMini0139.IsNearestHundredthReadout}
  reported is a nearest-hundredth readout of exact: it is an integer number of hundredths and differs from the exact value by at most half a hundredth
\end{definition}
\begin{proof}
  This is the declared model definition of Is Nearest Hundredth Readout.
\end{proof}
\begin{lemma}[signed Image Distance In Meters eq neg forty ninths]
  \label{lem:physics:phyx-mini-0139:phyxminiproblems-problemphyxmini0139-signedimagedistanceinmeters-eq-neg-forty-ninths}
  \lean{PhyXMiniProblems.ProblemPhyXMini0139.signedImageDistanceInMeters_eq_neg_forty_ninths}
  \uses{def:physics:phyx-mini-0139:phyxminiproblems-problemphyxmini0139-lengthinmeters, def:physics:phyx-mini-0139:phyxminiproblems-problemphyxmini0139-convexrearviewmirrorsetup, def:physics:phyx-mini-0139:phyxminiproblems-problemphyxmini0139-matchesproblemstatement, def:physics:phyx-mini-0139:phyxminiproblems-problemphyxmini0139-usescartesianconvexmirrorconvention, def:physics:phyx-mini-0139:phyxminiproblems-problemphyxmini0139-satisfiesconvexmirrorfocallaw, def:physics:phyx-mini-0139:phyxminiproblems-problemphyxmini0139-satisfiesgaussianmirrorequation}
  The curvature law and R = 16 m give f = -8 m; the Gaussian mirror equation with dₒ = 10 m then gives the virtual-image distance dᵢ = -40/9 m
\end{lemma}
\begin{proof}
  The conclusion follows from the governing relations and figure readouts named in the statement of signed Image Distance In Meters eq neg forty ninths.
\end{proof}
\begin{lemma}[transverse Magnification eq four ninths]
  \label{lem:physics:phyx-mini-0139:phyxminiproblems-problemphyxmini0139-transversemagnification-eq-four-ninths}
  \lean{PhyXMiniProblems.ProblemPhyXMini0139.transverseMagnification_eq_four_ninths}
  \uses{def:physics:phyx-mini-0139:phyxminiproblems-problemphyxmini0139-convexrearviewmirrorsetup, def:physics:phyx-mini-0139:phyxminiproblems-problemphyxmini0139-matchesproblemstatement, def:physics:phyx-mini-0139:phyxminiproblems-problemphyxmini0139-usescartesianconvexmirrorconvention, def:physics:phyx-mini-0139:phyxminiproblems-problemphyxmini0139-satisfiesconvexmirrorfocallaw, def:physics:phyx-mini-0139:phyxminiproblems-problemphyxmini0139-satisfiesgaussianmirrorequation, def:physics:phyx-mini-0139:phyxminiproblems-problemphyxmini0139-satisfiestransversemagnificationlaw}
  Substituting the derived virtual-image distance into m = -dᵢ/dₒ yields the exact positive, upright, reduced magnification m = 4/9
\end{lemma}
\begin{proof}
  The conclusion follows from the governing relations and figure readouts named in the statement of transverse Magnification eq four ninths.
\end{proof}
% --- Archon declaration topology end ---
% --- Archon physics formalization source end ---
